\subsection{Descriptions of Nucleon Structure}
\subsubsection{Parton Distribution Functions}
The PDF of a given hadron parameterizes the number density of a given type of parton at a given Bjorken-$x$, as formally introduced in Section~\ref{sec:factorization}.
All the possible partons are the 6 flavors of quarks and the one gluon, so for a given hadron there are 13 different parton distributions (counting anti-quarks).
For heavy-flavor quarks ($c$, $b$, $t$), whose masses are well above $\Lambda_{\mathrm{QCD}}$, their PDFs can be computed directly using pQCD.
For the light flavor quarks and gluon, the PDFs themselves must be determined from experimental data.

From DIS data, it is possible to obtain structure functions that, at leading order (LO), are sensitive to the quark content of the hadron across a wide range of $x$ and $Q^2$.
At LO, the structure functions are related to the quark PDFs according to Equation~\ref{eq:dis_to_pdf}~\cite{DISref}.
\begin{equation}
    \label{eq:dis_to_pdf}
    F_2(x, Q^2) = \frac{1}{2} \sum_{q} e^2_q \left[xf_q(x,Q^2) +xf_{\bar{q}}(x,Q^2)\right]
\end{equation}
where $e_q$ is the fractional charge of a given quark flavor.
The PDFs also have several sum rules that offer additional constraints to any fit~\cite{factorthrm}.
For illustrative purposes, all the following sum rules will apply to protons.
\begin{equation}
    \label{eq:mom_charge_sum_rule}
    \sum_{i}\int_{0}^{1} x f_i(x,Q^2) dx = 1,  \qquad \sum_{i}\int_{0}^{1} e_q f_i(x,Q^2) dx = 1
\end{equation}
Equation~\ref{eq:mom_charge_sum_rule} displays the sum rule conserving the total proton momentum (left) and the sum rule conserving the total proton charge (+1) (right).
\begin{equation}
    \label{eq:flavor_sum_rule}
    \int_{0}^{1} (f_u-f_{\bar{u}} ) dx = 2,  \;  \int_{0}^{1} (f_d-f_{\bar{d}} ) dx = 1,  \;\int_{0}^{1} (f_{\ge s}-f_{\ge\bar{s}} ) dx = 0, \;
\end{equation}
Equation~\ref{eq:flavor_sum_rule} shows the flavor-conservation sum rules for a proton with two $u$ valence quarks, one $d$ valence quark, and no strange or heavier valence quarks.
Under all of these constraints, dedicated PDF fitting groups will use fits of the form given in Equation~\ref{eq:pdf_fit}~\cite{pdg}.
\begin{equation}
    \label{eq:pdf_fit}
    f_i = x^a \left(\cdots\right)\left(1-x\right)^b
\end{equation}
In total the PDF fitting groups typically have approximately 14-32 free parameters in their fits.
The fits are optimized to minimize $\chi^2$ with respect to the input data, and the errors are determined using the Hessian method~\cite{pdg}.
For the light flavor quarks and gluon, while pQCD \textit{cannot} directly compute PDFs it \textit{can} describe how they should evolve with $Q^2$ through the Dokshitzer–Gribov–Lipatov–Altarelli–Parisi (DGLAP) evolution equations~\cite{DGLAP}.
The DGLAP equations evolve a prediction for a PDF at a scale $\mu_0$ to a different scale $\mu$.
PDF fitting groups generally start their PDFs at a scale of $\mu_0^2 \approx 1$--$4~\GeV^2$ and evolve them from there.
The evolution of the parton distribution functions with the factorization scale is illustrated in Figure~\ref{fig:herapdf_scales} for two representative values of $Q^2$~\cite{HERAF2}. The PDF fits were extracted using data from Figure~\ref{fig:proton_structure_function_f2}, among other sources.
\begin{figure}[H]
    \centering
    \begin{subfigure}[t]{0.49\linewidth}
        \centering
        \includegraphics[
            width=\linewidth,
            alt={Parton distribution functions at a factorization scale of 10 GeV squared.}
        ]{chapters/background/sections/nuclear_structure/figures/herapdf.pdf}
        \caption{$Q^2 = 10~\GeV^2$}
    \end{subfigure}
    \hfill
    \begin{subfigure}[t]{0.49\linewidth}
        \centering
        \includegraphics[
            width=\linewidth,
            alt={Parton distribution functions at a factorization scale of 10 TeV squared.}
        ]{chapters/background/sections/nuclear_structure/figures/herapdfhiq2.pdf}
        \caption{$Q^2 = 10~\TeV^2$}
    \end{subfigure}
    \caption{Proton PDFs at two factorization scales, showing the scale evolution predicted by the HERAPDF analysis. Figure from H1 and ZEUS Collaborations, arXiv:1506.06042v3~\cite{HERAF2}.}
    \label{fig:herapdf_scales}
\end{figure}

Similarly to DGLAP evolution, the Balitsky-Fadin-Kuraev-Lipatov (BFKL) equations describe the evolution of parton densities at small $x$~\cite{BFKL1, BFKL2}.

In general, many different processes and experiments contribute to global PDF analyses, and the neutral-current DIS example outlined above is merely illustrative.
The broad kinematic reach of the data used in global PDF analyses is shown in the left panel of Figure~\ref{fig:pdg_pdf_context}~\cite{pdg}, which includes data from fixed-target experiments, $e^-p$ collisions at HERA, and $pp$ ($p\bar{p}$) collisions at the LHC and Tevatron.

The right panel of Figure~\ref{fig:pdg_pdf_context} displays another example of a PDF fit.
Note that the valence and sea quarks are separate in the fit.
The valence quarks tend to carry a significant fraction of the parton's momentum and are peaked at $x=0.2$, while the contribution from sea quarks and gluons rises at low $x$.
In particular, this low-$x$ region is very interesting because of the behavior of the gluon PDF, which rises immensely due to soft-gluon radiation.
This rise at low $x$ is more prominent at high $Q^2$, where the higher-resolution probe can distinguish increasingly soft gluons.
By unitarity constraints, this gluon splitting cannot continue indefinitely as $x$ decreases.
At sufficiently small $x$, the gluon number density must decrease due to the recombination of these soft gluons in a phenomenon called \textit{gluon saturation}~\cite{gluonsat}.
\begin{figure}[H]
    \centering
    \begin{subfigure}[t]{0.49\linewidth}
        \centering
        \includegraphics[
            width=\linewidth,
            alt={Kinematic domains in Bjorken x and Q squared probed by fixed-target and collider experiments, including final states accessible at the LHC.}
        ]{chapters/background/sections/nuclear_structure/figures/kinematic_reach.pdf}
        \caption{Kinematic reach of fixed-target and collider experiments.}
    \end{subfigure}
    \hfill
    \begin{subfigure}[t]{0.49\linewidth}
        \centering
        \includegraphics[
            width=\linewidth,
            alt={NNLO MSHT20 global-analysis parton distribution functions at a scale of 10 squared GeV squared.}
        ]{chapters/background/sections/nuclear_structure/figures/nnlopdf.pdf}
        \caption{NNLO MSHT20 PDFs at $Q^2=10^2~\GeV^2$.}
    \end{subfigure}
    \caption{Examples of the experimental kinematic coverage and global PDF results. The right panel shows the parton distribution functions obtained from the NNLO MSHT20 global analysis at $Q^2=10^2~\GeV^2$. Figure from the Reference~\cite{pdg}.}
    \label{fig:pdg_pdf_context}
\end{figure}

\paragraph{Nuclear PDFs}
The partonic content of protons and neutrons (nucleons) is modified when they are bound into nuclei. This modification is typically quantified through a nuclear modification factor, which compares the parton distribution function in a nucleus with that in a free nucleon. The nuclear modification factor is defined as
\begin{equation}
    \label{eq:nuclear_modification_factor}
    R_i^A(x,Q^2) = \frac{f_i^A(x,Q^2)}{f_i(x,Q^2)},
\end{equation}
where $f_i^A$ is the PDF per nucleon in a nucleus with mass number $A$, and $f_i$ is the corresponding PDF in a free nucleon.  Values of $R_i^A$ below or above unity indicate suppression or enhancement, respectively, relative to the free-nucleon reference~\cite{nPDF}.
Figure~\ref{fig:npdf_modification} shows an illustration of the general shape of a nuclear modification factor, which can be decomposed into four main regions.
\begin{figure}[H]
    \centering
    \includegraphics[
        width=0.85\linewidth,
        alt={Illustration of the EPPS16 $R_i^A(x,Q^2)$ fit function.}
    ]{chapters/background/sections/nuclear_structure/figures/nPDFmod.pdf}
    \caption{Illustration of the EPPS16 $R_i^A(x,Q^2)$ fit function. Figure taken from Reference~\cite{nPDF}.}
    \label{fig:npdf_modification}
\end{figure}

\begin{enumerate}
    \item \textbf{Shadowing:} At low $x$, the density of very soft partons decreases in the nuclear environment because of a crowding effect from surrounding nucleons~\cite{shadowing}.
    \item \textbf{Antishadowing:} At more moderate $x$, the parton density is enhanced as partons redistributed from the shadowing region migrate to higher $x$~\cite{shadowing}.
    \item \textbf{EMC effect:} Measured by the European Muon Collaboration~\cite{EMCdiscovery}, this modification is not yet fully understood but is believed to be caused by short-range correlations between nucleon pairs~\cite{EMCmech}.
    \item \textbf{Fermi motion:} The quantum-mechanical motion of nucleons bound in the nucleus produces an enhancement at large $x$~\cite{nPDFFermi}.
\end{enumerate}

In general, nPDFs are less well determined than their free-proton counterparts, especially for nuclear species for which experimental data are sparse.

\FloatBarrier
\paragraph{Fragmentation Functions}
Fragmentation functions are the final-state analogues of PDFs and are written as $D_i^h(z,\mu_F^2)$.
They were introduced formally in Section~\ref{sec:factorization} as number-density functions representing the expected number of final-state hadrons carrying a momentum fraction $z$ of an outgoing parton.
Like PDFs, they are nonperturbative objects that must be determined from experiment.
A clean probe of fragmentation functions comes from the $e^+e^-\rightarrow h+X$ process, since there is no hadronic initial state and hence no PDF dependence.
The total cross section for this process can be used to determine the fragmentation functions in a manner similar to that used to determine structure functions in DIS.
The fragmentation functions can then be extracted from these measurements~\cite{pdg}.

\subsubsection{Generalized Parton Distributions and Color Transparency}

We have discussed already how PDFs encode the partonic content of hadrons as a function of the total \textit{longitudinal} momentum fraction of the hadron.
Further properties, such as the spatial distribution of partons, can also be of interest.

Electromagnetic form factors (FFs) describe the effective charge ($F_1(Q^2)$) and magnetic moment ($F_2(Q^2)$)\footnote{Despite the similar notation, these electromagnetic FFs differ from the structure functions introduced earlier.} of a hadron as a function of the momentum transfer $Q^2$.
For a proton, these factors are equivalent to the Mott-scattering limit as $Q^2\rightarrow0$, where $F_1(0)=1$ and $F_2(0)=\kappa_p$\footnote{$\kappa_p$ is the anomalous magnetic moment of the proton.}~\cite{EMFF}.
As $t$ increases, the FFs decrease as the scattering resolves further details of the charge and magnetic moments within the proton.

Next, we consider elastic scattering of a virtual photon off a proton.
Working in the light-front formalism, the momentum transfer ($Q^2$) lies in the transverse plane of the collision, so $Q^2 = -q^2_\perp$.
In this setting, a charge density can be extracted for the proton using the $F_1(Q^2)$ form factor through a Fourier transform from momentum to position space, as shown in Equation~\ref{eq:proton_density_FF}~\cite{Brodsky:CT}.
\begin{equation}
    \label{eq:proton_density_FF}
    \rho^p (\mathbf{b}_\perp) = \int \frac{d^2\mathbf{q}_\perp}{(2\pi)^2} \exp{\left(-i\mathbf{b}_\perp\cdot \mathbf{q}_\perp\right)} F(q^2)
\end{equation}

Deeply virtual Compton scattering (DVCS) is a type of deep ($Q^2\gg M^2$) exclusive scattering (DES) in which a photon or meson is electroproduced by scattering from a nucleon.
\begin{figure}[H]
    \centering
    \includegraphics[
        width=0.85\linewidth,
        alt={Handbag diagram of deeply virtual Compton scattering, showing a virtual photon scattering from a quark and a real photon produced from the nucleon.}
    ]{chapters/background/sections/nuclear_structure/figures/dvcshandbag.pdf}
    \caption{Handbag representation of deeply virtual Compton scattering, in which a virtual photon scatters from a parton in a nucleon and a real photon is produced. The process is described by generalized parton distributions~\cite{GPD:DVCS}.}
    \label{fig:dvcshandbag}
\end{figure}
As shown in Figure~\ref{fig:dvcshandbag}, this process has a factorized structure described by Generalized Parton Distributions (GPDs).
GPDs are powerful tools because they simultaneously access both the longitudinal momentum and transverse spatial distributions of partons within a hadron, unifying FFs and PDFs in one theoretical framework.
There are four GPDs for each quark and gluon, corresponding to the four independent helicity-spin combinations of the quark--nucleon system.
This discussion will focus on the GPD corresponding to an average over quark helicities and an unflipped nucleon spin, written as $H(x,\xi,t)$~\cite{GPD:Main}.
The three kinematic variables that parameterize GPDs are $x$, the average longitudinal momentum fraction of the struck quark (i.e., Bjorken $x$)\footnote{In the language of GPDs, $x$ can be negative. A negative $x$ corresponds to an antiquark with momentum fraction $|x|$.}; $x+\xi$ ($x-\xi$), the longitudinal momentum fraction of the struck quark before (after) the scattering; and $t$, the squared four-momentum transfer given to the nucleon~\cite{GPD:DVCS}.
The quark PDFs and FFs can both be derived from the GPDs through the relations given in Equation~\ref{eq:GPD_rules}~\cite{GPD:Main,Brodsky:CT}
\begin{equation}
    \label{eq:GPD_rules}
    H^q(x,0,0) =
    \begin{cases}
        f_q(x) & x > 0 \\
        -f_q(-x)  & x < 0
    \end{cases}
    \qquad \int_{-1}^{1} dx H^q(x,0,t) = F_1^q(t)
\end{equation}
Experimentally, GPDs can be determined by measuring DVCS rates, such as those measured at Jefferson Lab~\cite{GPD:DVCS}.
GPDs provide a richer description of nucleon structure than PDFs or FFs alone and can be used to understand parton dynamics in hadron collisions more closely.

\textit{Color transparency} (CT) is the tendency for a hadron in a configuration with a very high-$x$ parton to have a spatially compact transverse density~\cite{Brodsky:CT}.
The remnants of the struck hadron can fluctuate into a reduced-size, color-neutral configuration~\cite{New:CT}.
CT can be considered a QCD analogue of the charge-transparency phenomenon in QED~\cite{QED:CT}, but this comparison is complicated by the non-Abelian nature of QCD.
The dynamics underlying CT are closely related to GPDs, because GPDs encode both the transverse spatial density and the longitudinal momentum configurations of partons.

Consider a scattering with one struck parton in a proton and $n-1$ spectator partons.
The transverse coordinate $\mathbf{a}_\perp = \sum_{j=1}^{n-1} x_j \mathbf{b}_{\perp,j}$ represents the $x$-weighted transverse position of all the spectator partons.
Similarly to Equation~\ref{eq:proton_density_FF}, a Fourier transform of the GPD gives direct access to a \textit{transverse-distance-dependent PDF}, $f_q\left(x,\mathbf{a}_\perp\right)$, shown in Equation~\ref{eq:b_dep_PDF}.
\begin{equation}
    \label{eq:b_dep_PDF}
    f_q\left(x,\mathbf{a}_\perp\right) = \int \frac{d^2\mathbf{q}_\perp}{\left(2\pi\right)^2} \exp{\left(-i \mathbf{a}_\perp\cdot\mathbf{q}_\perp\right)} H^q(x,0,\mathbf{q}_\perp)
\end{equation}
Using this distribution, it is possible to derive the squared transverse spatial distance of the spectator partons as a function of the $x$ of the struck parton using Equation~\ref{eq:aT_xdep}~\cite{Brodsky:CT}.
\begin{equation}
    \label{eq:aT_xdep}
    \langle\mathbf{a}_\perp^2(x)\rangle_q = \frac{\int d^2\mathbf{a}_\perp \mathbf{a}_\perp^2 f_q\left(x,\mathbf{a}_\perp\right)}{\int  d^2\mathbf{a}_\perp f_q\left(x,\mathbf{a}_\perp\right)}
\end{equation}
CT effects become relevant as the $x$ of the struck quark approaches one. According to the PDFs, the high-$x$ struck quark is very likely to be a valence quark.
Therefore, considering only the valence quarks, the transverse radius of the proton can be written as a charge-weighted sum, $\langle\mathbf{a}_\perp^2(x)\rangle_p = 2 e_u \langle\mathbf{a}_\perp^2(x)\rangle_u + e_d \langle\mathbf{a}_\perp^2(x)\rangle_d$.
Figure~\ref{fig:ct_transverse_size} shows a prediction for $\langle\mathbf{a}_\perp^2(x)\rangle_p$ that assumes a phenomenological functional form for how the GPD encodes the spatial distribution of the spectator quarks as a function of $x$. The prediction is compared with data from the CLAS and HERMES experiments and shows good agreement within the experimental uncertainties.
\begin{figure}[H]
    \noindent
    \begin{subfigure}[t]{0.48\textwidth}
        \centering
        \includegraphics[
            width=\linewidth,
            alt={Proton transverse size squared as a function of the struck parton's momentum fraction, comparing a model prediction with CLAS and HERMES data.}
        ]{chapters/background/sections/nuclear_structure/figures/CT1D.pdf}
        \caption{Model prediction compared with CLAS and HERMES data.}
    \end{subfigure}%
    \hfill%
    \begin{subfigure}[t]{0.48\textwidth}
        \centering
        \includegraphics[
            width=\linewidth,
            alt={Illustration of the proton transverse size decreasing at large light-front momentum fraction, showing the transition from infrared to ultraviolet behavior.}
        ]{chapters/background/sections/nuclear_structure/figures/CT2D.pdf}
        \caption{Color-transparency behavior at large light-front momentum fraction.}
    \end{subfigure}
    \caption{The $x$ dependence of the proton transverse size, $\langle\mathbf{a}_\perp^2(x)\rangle_p$, determined from the hadron profile function. The left panel compares the model with CLAS and HERMES data. The right panel shows an illistraion of the CT phenomenon. At large $x$, the transverse size approaches that of a point-like color-singlet configuration. Figure taken from Reference~\cite{Brodsky:CT}.}
    \label{fig:ct_transverse_size}
\end{figure}

The CT phenomenon has also been theoretically framed in terms of color fluctuations (CF)~\cite{CF1,CF2,CF3}, which parameterize the reduced interaction strength of the struck-hadron remnant after the high-$x$ parton has been removed.
Whichever term is preferred, this phenomenon is particularly relevant for proton--nucleus collisions.
When the projectile proton undergoes a hard scattering and enters a spatially compact state, the probability of interacting with other nucleons in the nucleus decreases.
This is schematically shown in Figure~\ref{fig:color_fluctuation_pA}.
\begin{figure}[H]
    \centering
    \includegraphics[
        width=0.85\linewidth,
        alt={Schematic proton--nucleus collision with a fixed target geometry. A weakly interacting compact proton and a more strongly interacting proton are shown, with impacted nucleons highlighted in red.}
    ]{chapters/background/sections/nuclear_structure/figures/pACT.pdf}
    \caption{The same target geometry is shown for a compact proton configuration with a smaller interaction probability and for a larger configuration that interacts with more nucleons. The red tube indicates the projectile's transverse extent through the nucleus, while red nucleons mark interactions.  Taken from Reference~\cite{CF3}.}
    \label{fig:color_fluctuation_pA}
\end{figure}
